\section{Minimal-polynomial profiles for actual Heisenberg sources}
\label{sec:heisenberg-minimal-polynomial}

\paragraph{Status and provenance.}
D-TD61, question Q2fq, is a new all-prime refinement of
the short projected-tuple boundary.
The parent proposed using the orbit span of the low-rank generators
under one remaining generator, with label irreducibility explicit.
The author derived the cyclic-module orbit-count trace estimate
below and its quantitative combination with that span argument.
The complete written proof, code and output have passed the author's
whole reading; the exact supplementary certificate passed its first run.
Complete root and independent whole-proof reviews also pass.
No full Heisenberg exponent bound is asserted.

\paragraph{Actual factor and all original periods.}
Let an actual connected finite $J$-Galois sphere cover have a faithful
irreducible complex coefficient representation on $V$, with actual
generating product-one inertia $T_1,\ldots,T_r$ satisfying
\[
 e=\dim V,\qquad
          \sum_{i=1}^r\operatorname{rank}(T_i-I)=2e+6 .
 \tag{HMP.1}
\]
Suppose $V=W\otimes U$ is preserved factorwise and its actual finite
projective image $P$ on $W$ contains a normal nondegenerate projective
Heisenberg group
\[
 Q\simeq\mathbf F_p^{\,2m},\qquad m\ge1,\quad
                      a=\dim W=p^m ,
 \tag{HMP.2}
\]
acting irreducibly on $W$.  The prime $p$ is arbitrary, including two.
Let $y_i\in P$ be the projected inertia and let
$g_i\in\operatorname{Sp}(Q)$ be its action on the label space.
After omitting the identity $y_i$ and relabelling, there are $k$
projected entries.  Their label images still have product one and
generate the actual group $H\subset\operatorname{Sp}(Q)$.
Some $g_i$ can be the identity even when $y_i\ne1$.
For a label operator $g$ define
\[
                    D(g)=\deg\operatorname{minpoly}_{\mathbf F_p}(g).
 \tag{HMP.3}
\]
There is no twelve-mark, whole-original support-floor, or self-duality
premise on the coefficient representation.  The whole partner $U$
is arbitrary.  Primitivity of $P$ is not required in the first two
conclusions below; irreducibility of $H$ is a separate hypothesis
for the final conclusion.

\begin{theorem}[D-TD61]
Let $A$ be a finite-projective-order normalizer of the represented
group $Q$, and let $g$ be its label action.  For every integer
$R$ with $0\le R<D=D(g)$, the all-target bound
\[
 \frac{\nu_W(A)}a\ \ge\
  1-\left(\frac1{R+1}+\frac43p^{\,2R-D}\right)^{1/2}
 \tag{HMP.4}
\]
holds, with no restriction on the full original period used to
normalize $A$.  A negative right side merely gives a vacuous bound.
It survives tensoring with any entire finite-order partner.
In particular,
\[
               D(g)\ge160
          \quad\Longrightarrow\quad
                    \frac{\nu_W(A)}a>\frac{22}{25}.
 \tag{HMP.5}
\]

For the actual source (HMP.1)--(HMP.2), at most one of its
nonidentity projected entries has $D(g_i)\ge160$.
If, additionally, the actual $\mathbf F_pH$-module $Q$ is
irreducible, then for every two distinct indices $i,j$ among
the $k$ projected entries,
\[
 2m\ \le\
 159\max\left\{1,\ \sum_{\ell\ne i,j}
                             \operatorname{rank}_{\mathbf F_p}(g_\ell-I)
       \right\}.
 \tag{HMP.6}
\]
Consequently, if all but two label ranks are at most $R_0\ge0$, then
\[
          2m\le159\max\{1,(k-2)R_0\}.
 \tag{HMP.7}
\]
For $k=3,4$ and $R_0=10$, this gives respectively $m\le795$
and $m\le1590$.  It does not remove irreducibility of the label
module from this last assertion.
\end{theorem}

\begin{proof}
\emph{Step 1: a cyclic invariant subspace of the exact degree.}
Put $f=\operatorname{minpoly}(g)$ and $D=\deg f$.
There is a $g$-invariant cyclic subspace $C\subset Q$ of dimension
$D$, on which the minimal polynomial is $f$.
Here is an elementary justification over the original field.
Factor $f=\prod_\alpha f_\alpha^{b_\alpha}$ into powers of distinct
monic irreducibles.  Bezout identities give the direct primary
decomposition of $Q$.  In its $f_\alpha$-primary summand choose
$v_\alpha$ not killed by $f_\alpha(g)^{b_\alpha-1}$.
Such a vector exists by minimality of the exponent.
Its annihilator is exactly $f_\alpha^{b_\alpha}$.
Thus $v=\sum_\alpha v_\alpha$ has annihilator $f$ and
\[
           C=\mathbf F_p[g]v\simeq \mathbf F_p[x]/(f),
                         \qquad \dim C=D .
 \tag{HMP.8}
\]
The invertibility of $g$ makes the constant term of $f$ nonzero.

For a monic polynomial $h$ of degree $j$, multiplication by $h$
on this cyclic module has kernel dimension
$\deg\gcd(f,h)\le j$.  Indeed, writing $f=dc$, $h=dh'$ with
$d=\gcd(f,h)$ and $\gcd(c,h')=1$, the condition $f\mid hz$ is
equivalent to $c\mid z$; that ideal modulo $f$ has dimension
$\deg d$.  In particular at most $p^j$ vectors are killed by $h$.
There are exactly $p^j$ monic polynomials of degree $j$.
It follows by a union bound, including the zero vector at $j=0$,
that the proportion of vectors in $C$ whose cyclic span has
dimension at most $R$ is at most
\[
       p^{-D}\sum_{j=0}^R p^{2j}
           \le \frac43p^{\,2R-D}.
 \tag{HMP.9}
\]
The last inequality uses $p\ge2$.  Repeated factors of $f$ cause
no exception; no semisimplicity is assumed.

\emph{Step 2: cyclic orbits and the correct fixed-rank comparison.}
Let $L$ be any positive integer with $g^L=I$.
The number of $\langle g\rangle$-orbits on $C$, divided by $|C|$,
is the average over $v\in C$ of $1/|\langle g\rangle v|$.
A vector whose cyclic span has dimension greater than $R$
has orbit length at least $R+1$, since that span is generated
by the vectors in its orbit.  Equation (HMP.9) therefore gives
\[
 \frac{\#(C/\langle g\rangle)}{|C|}
       \le\frac1{R+1}+\frac43p^{\,2R-D}.
 \tag{HMP.10}
\]
Burnside's elementary orbit count, applied also to a nonfaithful
cyclic action of period $L$, identifies the left side with
$L^{-1}\sum_{j=0}^{L-1}p^{-\operatorname{rank}_C(g^j-I)}$.
For completeness, each orbit contributes exactly $L$ pairs
$(v,j)$ with $g^jv=v$, which proves that identity by double counting.

Because $C$ is an invariant subspace of $Q$, the image on $C$
is a subspace of the image on $Q$.  Thus the comparison has
the direction
\[
 \operatorname{rank}_Q(g^j-I)\ge\operatorname{rank}_C(g^j-I),
 \qquad
 p^{-\operatorname{rank}_Q(g^j-I)}
           \le p^{-\operatorname{rank}_C(g^j-I)}.
 \tag{HMP.11}
\]
Neither the complement of $C$ nor an invariant symplectic
complement is needed.

\emph{Step 3: the entire original-period Fourier estimate.}
The finite-projective-order operator $A$ and the finite normal
subgroup $Q$ generate a finite projective group.  Its inverse image
in $\operatorname{SL}(W)$ is finite.  A common invariant positive
Hermitian form therefore exists.
Choose any $L$ with $A^L$ scalar; then $g^L=I$, so Step~2 applies.
Rescale $A$ so that $A^L=I$ and use this invariant form.
An arbitrary period of the label $g$ alone is not substituted here.
The phase-sensitive trace formula of TD05 gives for every power
\[
       \frac{|\operatorname{tr}A^j|^2}{a^2}
                \le p^{-\operatorname{rank}_Q(g^j-I)} .
 \tag{HMP.12}
\]
This includes powers with trivial label or projective action:
their trace can be zero or can have its full scalar magnitude.
Cauchy--Schwarz, (HMP.10)--(HMP.12), and Fourier averaging at an
arbitrary target eigenvalue $\lambda$ now give
\[
 \begin{split}
 \frac{\dim\ker(A-\lambda I)}a
 &\le \frac1L\sum_{j=0}^{L-1}\frac{|\operatorname{tr}A^j|}{a}\\
 &\le\left(\frac1L\sum_{j=0}^{L-1}
                    \frac{|\operatorname{tr}A^j|^2}{a^2}\right)^{1/2}\\
 &\le\left(\frac1{R+1}+\frac43p^{\,2R-D}\right)^{1/2}.
 \end{split}
 \tag{HMP.13}
\]
Scaling back does not change the largest eigenspace.

For an actual coefficient inertia $T_i=A_i\otimes B_i$ of order
$L_i$, the two $L_i$-th powers are reciprocal scalars.
Reciprocal rescaling makes both powers identity without changing
$T_i$.  Thus (HMP.13) uses the full original period $L_i$;
all repeated label periods and scalar powers are retained.
Diagonalizing the entire $B_i$ shows that tensor convolution
cannot increase the largest relative eigenspace beyond that of
$A_i$.  This proves (HMP.4) and its arbitrary-partner assertion.

\emph{Step 4: an absolute degree threshold and the original budget.}
If $D\ge160$, take $R=70$ and use $p\ge2$.  The exact bound is
\[
 \frac1{71}+\frac43p^{140-D}
   \le\frac1{71}+\frac4{3\cdot2^{20}}
   <\frac9{625}.
 \tag{HMP.14}
\]
Indeed $9/625-1/71=14/44375$, so the strict comparison is
$4\cdot44375<42\cdot2^{20}$.
This proves (HMP.5).

The projected tuple contains at least three nonidentity entries.
Otherwise product one and generation make $P$ cyclic, contrary
to its noncyclic subgroup $Q$.
Suppose two of these entries have $D(g_i)\ge160$.
They each cost more than $22e/25$ in (HMP.1).
A third nonidentity projected entry costs at least $e/4$
by the all-prime Heisenberg floor of TD08.
Consequently
\[
          2+\frac6e>\frac{44}{25}+\frac14
                           =\frac{201}{100},
                \qquad e<600.
 \tag{HMP.15}
\]
On the other hand $160\le D(g_i)\le2m$, so
$e\ge p^m\ge2^{80}>600$.  This contradiction proves that at most
one such entry exists.  No cost was assigned to any omitted mark.

\emph{Step 5: the irreducible-label orbit span.}
Assume now that $Q$ is irreducible as an $\mathbf F_pH$-module.
Fix distinct indices $i,j$.  At least one of their minimal-polynomial
degrees is at most $159$; interchange the two names so that it
is $g=g_i$.  The product-one relation expresses $g_j$ in the
subgroup generated by $g_i$ and the $g_\ell$ with $\ell\ne i,j$.
Therefore
\[
                 H=\langle g,\ g_\ell\ (\ell\ne i,j)\rangle .
 \tag{HMP.16}
\]
Let $R_{ij}=\sum_{\ell\ne i,j}\operatorname{rank}(g_\ell-I)$.
Consider the actual subspace
\[
 X=\sum_{\ell\ne i,j}\ \sum_{b\in\mathbf Z}
                         g^b\operatorname{im}(g_\ell-I)
          \ \subset Q .
 \tag{HMP.17}
\]
It is $g$-stable.  Each $g_\ell$ preserves $X$, since
$(g_\ell-I)X\subset\operatorname{im}(g_\ell-I)\subset X$.
Invertibility gives invariance rather than just containment.
Thus (HMP.16) makes $X$ an $H$-invariant subspace.
The nonzero constant term of the minimal polynomial expresses
all positive and negative powers of $g$ in the span of
$I,g,\ldots,g^{D(g)-1}$.  Hence
\[
                          \dim X\le D(g)R_{ij}.
 \tag{HMP.18}
\]
If $R_{ij}>0$, then $X\ne0$ and label irreducibility gives $X=Q$,
so $2m\le159R_{ij}$.
If $R_{ij}=0$, all the other label generators are identity, so
$H=\langle g\rangle$.  The cyclic span of any nonzero vector is
then all of $Q$ by irreducibility, giving $2m=D(g)\le159$.
This proves (HMP.6), and (HMP.7) follows immediately.
\end{proof}

\paragraph{When the additional label hypothesis is available.}
If $Q$ is minimal normal in its actual projective image $P$, its
nonzero proper $\mathbf F_p$-subspaces cannot be $H$-invariant,
since they would be normal subgroups of $P$.  Thus the required
label irreducibility holds in that case.
At a Heisenberg step of the complete TD12 recursion, the chosen
minimal normal elementary abelian subgroup maps faithfully onto
$Q$ in the extracted coarse projective quotient.  An invariant
proper subspace there would pull back to an invariant subgroup
of the chosen minimal normal subgroup.  This supplies the same
irreducibility at those steps.
It is not inferred from primitivity of an arbitrary Heisenberg
normalizer alone.

\paragraph{Strict short-tuple consequence and boundary.}
In the binary case, TD59 separately gives $k=3,4$ when $m>80$
and the selected projective image is primitive.
If its label action is also irreducible, the present theorem
excludes $m>795$ with $k=3$ and any label rank at most ten,
and excludes $m>1590$ with $k=4$ and two label ranks at most ten.
It leaves all other short profiles open.
The all-prime at-most-one conclusion itself needs neither
primitivity nor label irreducibility.
Nothing here bounds all Heisenberg exponents, constructs a source,
or resolves original-field, natural-polarization or detecting-cycle
questions.  No unproved bounded-order signature consequence is included.
